\section{Paradigma Divide and Conquer}
\subsection{Proprietà di sottostruttura}
\subsubsection*{Idea generale}
Il paradigma divide and conquer (divide et impera) è un approccio che consiste nel suddividere un'istanza di grandi dimensioni
in una o più sottoistanze di dimensioni inferiori, più semplici da risolvere, per poi combinare le soluzioni di queste istanze ed
ottenere la soluzione dell'istanza originale.

\subsubsection*{Proprietà di sottostruttura}
Dato un problema \(\pi \subseteq \mathcal{I} \times \mathcal{S}\), il paradigma si compone di tre componenti:
\begin{itemize}[topsep=5pt, itemsep=3pt] 
	\item \(A_D : \mathcal{I} \to \mathcal{I}*\), funzione di divide che suddivide un'istanza \(i \in \mathcal{I}\) di
	taglia \(|i| > n_0\) in una sequenza di una o più sottoistanze più piccole \(\langle i_1, i_2, \dots i_k \rangle\)
	di taglia \(|i_j| < |i|\) per ogni \(j\)
	\item \(A_C : \mathcal{S}* \to \mathcal{S}\), funzione di combine che combina la sequenza di soluzioni
	\(\langle s_1, s_2, \dots s_k \rangle\) associate alle sottoistanze \(\langle i_1, i_2, \dots i_k \rangle\)
	nella soluzione \(s\) associata all'istanza \(i\)
	\item \(D\&C : \mathcal{I} \to \mathcal{S}\), chiamata ricorsiva che trasforma individualmente ogni sottoistanza \(i_j\)
	nella corrispettiva soluzione \(s_j\), eventualmente applicando le due funzioni \(A_D\) e \(A_C\) in modo ricorsivo
\end{itemize}

\noindent
Le due funzioni \(A_D\) e \(A_C\) devono soddisfare la seguente proprietà di sottostruttura:
\[\forall i \in \mathcal{I}, |i| > n_0 \text{ vale } \left[ \; D(i) = \; \langle i_1, \dots, i_k \rangle \; \wedge \; \forall j \; (i_j, s_j) \in \pi \; \wedge \; C(\langle s_1, \dots, s_k \rangle) = s \implies (i, s) \in \pi \; \right]\]

\subsubsection*{Chiamata ricorsiva (conquer) e definizione del caso base}
La chiamata ricorsiva \(D\&C(i)\) su un'istanza \(i\) è definita come segue:
\begin{itemize}
	\item se \(|i| > n_0\) allora \(D\&C(i) = A_C(D\&C^*(A_D(i)))\), ovvero riapplica il modello di divide and conquer
	\item se \(|i| \leq n_0\) allora utilizza una soluzione diretta ed immediata per risolvere \(i\)
\end{itemize}
In generale il parametro \(n_0\) è scelto in base al problema. Può essere il minimo valore per cui è possibile applicare
il modello di divide and conquer, oppure un valore sufficientemente grande sotto di cui si conosce un algoritmo più
efficiente per risolvere il problema.

\begin{algo}{D\&C}{i}{\text{istanza \(i \in \mathcal{I}\)}}{\text{soluzione \(s \in \mathcal{S}\)}}
\If{\(|i| \leq n_0\)} \Comment{caso base \(\qquad\qquad\qquad\qquad\qquad\qquad\qquad\qquad\qquad\qquad\;\)}
	\State \(s \gets A_\text{diretto}(i)\) \Comment{soluzione diretta ed immediata \(\quad\qquad\qquad\qquad\quad\,\)}
	\State \textbf{return} \(s\)
\Else \Comment{passo ricorsivo \(\qquad\qquad\qquad\qquad\qquad\qquad\qquad\qquad\qquad\)}
	\State \(\langle i_1, i_2, \dots i_k\rangle \; \gets A_D(i)\) \Comment{divide \(\quad\qquad\qquad\qquad\qquad\qquad\qquad\qquad\qquad\qquad\)}
	\For{\(j \gets 1\)}{\(k\)} \Comment{conquer \(\quad\qquad\qquad\qquad\qquad\qquad\qquad\qquad\qquad\quad\:\)}
		\State \(s_j \gets D\&C(i_j)\)
	\EndFor
	\State \(s \gets A_C(\langle s_1, s_2, \dots s_k \rangle)\) \Comment{combine \(\quad\qquad\qquad\qquad\qquad\qquad\qquad\qquad\qquad\quad\)}
	\State \textbf{return} \(s\)
\EndIf
\end{algo}

\newpage

\subsection{Analisi della correttezza}
\subsubsection*{Proprietà di sottostruttura in funzione di un parametro intero}
Per dimostrare la correttezza di un algoritmo divide and conquer, si utilizza la proprietà di sottostruttura e il principio
di induzione in forma forte. Per prima cosa si esprime la proprietà di sottostruttura in funzione di un parametro intero,
in modo da renderla più adatta alla dimostrazione induttiva, ovvero:
\[\forall i \in \mathcal{I} \; \underbrace{\, \left[ \; (i, \; D\&C(i)) \in \pi \; \right] \,}_{P(i)} \;\; = \;\; \forall n \in \mathbb{N} \; \underbrace{\, \left[ \forall i \in \mathcal{I} \text{ con } |i| = n \text{ vale } (i, D\&C(i)) \in \pi \; \right] \,}_{P(n)}\]

\subsubsection*{Dimostrazione con il principio di induzione in forma forte}
Si dimostra, infine, che \(P(n)\) è vera per ogni \(n \in \mathbb{N}\) tramite il principio di induzione in forma forte:
\begin{itemize}
	\item si dimostra che \(P(n)\) è vera per ogni \(n \leq n_0\) (caso base)
	\item si dimostra che se \(P(n)\) è vera per tutti i valori \(n \leq m\) allora anche \(P(m + 1)\) è vera (passo induttivo)
\end{itemize}

\subsection{Albero della ricorsione}
\subsubsection*{Albero della ricorsione}
L'albero della ricorsione è una rappresentazione grafica che delle chiamate ricorsive di un algoritmo divide and
conquer con le seguenti caratteristiche:
\begin{itemize}
	\item ogni nodo rappresenta una chiamata ricorsiva \(D\&C(i)\) per un'istanza \(i \in \mathcal{I}\)
	\item ogni nodo ha come figli i nodi che rappresentano le chiamate ricorsive \(D\&C(i_j)\) per ogni sottoistanza
	\(i_j\) ottenuta dalla procedura di divide \(A_D(i)\)
	\item la taglia dell'istanza dei nodi figli è inferiore alla taglia dell'istanza del nodo padre, in quanto la taglia
	dell'istanza diminuisce progressivamente ad ogni chiamata ricorsiva
	\item ogni foglia rappresenta una chiamata ricorsiva \(D\&C(i)\) per un'istanza \(i \in \mathcal{I}\) di taglia
	\(|i| \leq n_0\), ovvero un caso base
\end{itemize}

\subsubsection*{Esplorazione dell'albero della ricorsione}
L'esplorazione dell'albero della ricorsione avviene secondo la visita in-order, ovvero prima si visita il figlio sinistro,
poi il nodo padre, poi il figlio destro. In questo modo, dato il nodo dell'albero che si sta attualmente visitando, si è
certi che tutti i nodi a sinistra di esso sono stati visitati, e i nodi a destra di esso sono ancora da visitare.

La costruzione della soluzione avviene secondo la visita post-order, ovvero prima si risolvono le istanze dei figli, poi
viene costruita la soluzione per il nodo padre, secondo l'algoritmo DFS (Depth First Search).

Lo spazio occupato è proporzionale allo spazio richiesto da ciascuna chiamata ricorsiva e dal numero di chiamate ricorsive
attive contemporaneamente, ovvero alla profondità dell'albero della ricorsione.

\subsubsection*{Cenni sulla ricorsione nei primi linguaggi di programmazione}
La memoria a disposizione sui primi computer era molto limitata e non era sufficiente per implementare uno stack che isolasse
le variabili locali delle varie chiamate ricorsive attive. I primi linguaggi di programmazione, infatti, salvavano le variabili
locali in posizioni fisse di memoria, comuni ad ogni chiamata ricorsiva. Per questo motivo, all'inizio la ricorsione non era
supportata. Il primo linguaggio di programmazione a supportare nativamente la ricorsione è stato il linguaggio LISP, sviluppato
da John McCarthy nel 1958, ma il suo utilizzo rimase fortemente scoraggiato per anni.

Ad oggi, la ricorsione è supportata nativamente da tutti i principali linguaggi di programmazione ed esistono ottimizzazioni
a livello hardware (istruzioni macchina per la gestione dello stack) per renderla efficiente e non sprecare memoria.

\subsection{Analisi della complessità e formulone}
\subsubsection*{Equazione alle ricorrenze}
L'equazione alle ricorrenze è un'equazione che esprime la complessità temporale di un algoritmo divide and conquer
in funzione della complessità temporale delle chiamate ricorsive e delle operazioni di divide e combine. Di seguito
è riportata la forma generale con le proprietà indicate sotto:
\[T(n) \; = \; \begin{cases}
	\; T_0(n) & \text{se } n \leq n_0 \quad - \quad \text{caso base} \\[3pt]
	\; s(n) \cdot T(f(n)) + \omega(n) & \text{se } n > n_0 \quad - \quad \text{passo ricorsivo}
\end{cases}\]
\begin{itemize}
	\item \(T_0(n)\) è il costo dell'algoritmo diretto per risolvere un'istanza di taglia \(n \leq n_0\)
	\item \(s(n)\) è il numero di chiamate ricorsive effettuate per risolvere un'istanza di taglia \(n\) e coincide con
	il numero di nodi figli nell'albero della ricorsione
	\item \(f(n)\) è la taglia della sottoistanza, generata dalla procedura di divide eseguita su un'istanza di taglia
	\(n\), che viene passata alla chiamata ricorsiva
	\item \(T(f(n))\) è il costo di ogni chiamata ricorsiva per risolvere un'istanza di taglia \(f(n)\)
	\item \(\omega(n) = T_D(n) + T_C(n)\) è il costo complessivo delle procedure di divide e di combine
\end{itemize}

\vspace{5pt}

\noindent
Per ottenere il valore analitico ed eliminare le ricorrenze è necessario risolvere l'equazione attraverso un'analisi
analitica, o ricorrendo al formulone o al master theorem.

\subsubsection*{Iterate di \(f\)}
Si definisce l'iterata di una funzione \(f\) come la potenza della composizione di \(f\) con se stessa.
\[\begin{cases}
	f^{(0)}(n) = n \\[3pt]
	f^{(l)}(n) = f(f^{(l-1)})(n) & \text{per } l > 0
\end{cases}\]
Si osserva che \(f^{(i)}(n)\) corrisponde alla taglia dell'istanza all'i-esima chiamata ricorsiva. In particolare,
inizialmente \(f^{(i)}(n) = n\) per \(i = 0\) e progressivamente diminuisce ad ogni iterata, fino a raggiungere il caso
base \(f^{(k+1)}(n) \leq n_0\) per un determinato valore di \(i = k\). Si definisce \(k\) come il massimo valore di \(i\)
per cui la taglia dell'istanza è ancora maggiore di \(n_0\), ovvero:
\[k = f^*(n, n_0) := \max \left\{ k : f^{(k)}(n) > n_0 \right\}\]

\subsubsection*{Formulone}
Il formulone è una formula generale che risolve l'equazione generale alle ricorrenze indicata sopra.
\[
T(n) = \underbrace{
	\overbrace{\sum_{l=0}^{f^* (n, \, n_0)}}^{\begin{array}{c} \text{\scriptsize per ogni} \\[-3pt] \text{\scriptsize livello interno} \end{array}}
	\overbrace{\left[ \prod_{j=0}^{l-1} s \left(f^{(j)}(n) \right) \right. \!}^{\begin{array}{c} \text{\scriptsize \# nodi al} \\[-3pt] \text{\scriptsize livello \(l\)-esimo} \end{array}}
	\; \cdot \!\!\! \overbrace{\left. \phantom{\prod_{j}^{l} \!\!\!\!\!\!\!\!} \omega \left( f^{(l)}(n) \right) \right]}^{\begin{array}{c} \text{\scriptsize costo interno} \\[-3pt] \text{\scriptsize di ogni nodo} \\[-3pt] \text{\scriptsize al livello \(l\)-esimo} \end{array}}
	}_\text{lavoro compiuto nei nodi interni}
	\;\; + \;\;\;
	\underbrace{ \overbrace{\prod_{j=0}^{f^*(n, n_0)} s \left(f^{(j)}(n) \right) \!}^{\begin{array}{c} \text{\scriptsize \# di foglie} \end{array}}
	\; \cdot \; \overbrace{ \phantom{\prod_{j}^{f^*} \!\!\!\!\!\!\!\!} T_0 \left( f^{(f^*(n, n_0) + 1)}(n) \right)}^{\begin{array}{c} \text{\scriptsize costo di ogni foglia}\end{array}}
	}_\text{lavoro compiuto nelle foglie}
\]

\vspace{7pt}

\noindent
Il formulone introduce le seguenti piccole semplificazioni che devono essere soddisfatte affinché la formula
indicata sopra sia valida:
\begin{itemize}
	\item il numero di figli dei nodi di un certo livello è costante, ovvero \(s(f^{(l)}(n)) = s_l\) per ogni livello \(l\)
	\item la taglia di tutti i nodi di un certo livello è costante, ovvero \(f^{(l)}(n) = f_l\) per ogni livello \(l\)
\end{itemize}

\subsubsection*{Formulone applicato ad un caso particolare}
Di seguito un esempio di applicazione del formulone ad un caso particolare con:
\begin{itemize}
	\item \(s(n) = a\): numero di figli costante per ogni nodo interno
	\item \(f(n) = n/2\): la taglia dell'istanza si dimezza ad ogni chiamata ricorsiva
	\item \(\omega(n) = \beta \, n\): costo di divide and conquer lineare per ogni nodo interno
	\item \(T_0(n) = 1\): costo costante per ogni foglia
\end{itemize}
\begin{align*}
	T(n) \;\; &= \sum_{l=0}^{(\log_2 n) - 1} \left( a^l \cdot \beta \; \frac{n}{2^l} \right) \; + \; a^{\log_2 n} \cdot T_0(n) \; = \; \beta \, n \cdot \! \sum_{l=0}^{(\log_2 n) - 1} \!\! \left( a/2 \right)^l \; + \; n^{\log_2 a} = \\[3pt]
	&= \;\; \beta \, n \cdot \frac{(a/2)^{\log_2 n} - 1}{(a/2) - 1} + n^{\log_2 a} \; = \; \frac{\beta}{a/2 - 1} \cdot n \left( n^{\log_2 a/2} - 1\right) + n^{\log_2 a} = \\[3pt]
	&= \; \underbrace{\frac{\beta}{a/2 - 1} \cdot (n^{\log_2 a} - n)}_{\text{costo dei nodi interni}} \;\; + \!\! \underbrace{ \;\; n^{\log_2 a} \;\;}_{\text{costo delle foglie}}
	\!\! = \;\; \Theta \left( n^{\log_2 a} \right)
\end{align*}

\subsubsection*{Complessità degli algoritmi di ordinamento basati sui confronti}
Al momento non si conoscono funzioni che legano problemi alla relativa complessità asintotica e più in generale non si
conoscono metodi per individuare in maniera assoluta i lower bound di problemi, ad eccezione di quelli deducibili banalmente.

Ad esempio gli algoritmi di ordinamento basati sui confronti hanno lower bound banale di \(\Theta(n \log n)\), siccome
\(n \log n\) è la quantità di informazione necessaria a rappresentare tutte le possibili permutazioni di \(n\) elementi ed
equivale al numero minimo di confronti (rappresentabili in bit) necessari per distinguere tutte le possibili permutazioni.
\[\# \; \text{permutazioni} = n! \approx n^n \qquad\quad \text{informazione}(\text{perm} \,_i) = \log_2(n!) \approx \log_2(n^n) = n \log_2 n\]

\subsection{Master theorem}
\subsubsection*{Assunzioni iniziali}
Il master theorem è un teorema che fornisce una soluzione asintotica immediata in funzione di due parametri \(a\) e \(b\)
per una classe di equazioni alle ricorrenze con le seguenti caratteristiche:
\begin{itemize}
	\item \(s(n) = a\): numero di figli costante per ogni nodo interno
	\item \(f(n) = n/b\): la taglia dell'istanza si divide per una costante \(b > 1\) ad ogni chiamata ricorsiva
	\item \(n_0 = 1\): caso base per istanze di taglia 1
\end{itemize}

\[T(n) \; = \; \begin{cases}
	\; T_0(n) & \text{se } n \leq 1 \\[3pt]
	\; a \cdot T(n/b) + \omega(n) & \text{se } n > 1
\end{cases}\]

\subsubsection*{Funzione soglia}
Si definisce una \textbf{funzione soglia} \(\sigma(n)\) che rappresenta il numero di foglie dell'albero della ricorsione.
Tale funzione permette di calcolare in maniera immediata il costo complessivo dovuto al lavoro compiuto nelle foglie che
funge da lower bound della complessità dell'algoritmo.
\[\sigma(n) = n^{\log_b a} \qquad T(n) \, \geq \, T_{\text{foglie}}(n) \, = \, T_0(n) \cdot \# \; \text{foglie} \, = \, T_0(n) \cdot \sigma(n)\]

\subsubsection*{Classi di funzioni e relative complessità asintotiche}
Paragonando il costo del lavoro compiuto da ogni nodo interno \(\omega(n)\) con la funzione soglia \(\sigma(n)\), si
individuano tre casi per cui è possibile calcolare la complessità asintotica di \(T(n)\) in maniera immediata:
\[T(n) \; = \; \begin{cases}
	\; \displaystyle \Theta (n^{\log_b a}) & \text{\footnotesize se il lavoro nei nodi interni \(\omega(n)\) è trascurabile rispetto a quello delle foglie \(\sigma(n)\)} \\[3pt]
	\; \displaystyle \Theta (n^{\log_b a} \, \log_b n) & \text{\footnotesize se il lavoro nei nodi interni \(\omega(n)\) è dello stesso ordine di quello delle foglie \(\sigma(n)\)} \\[3pt]
	\; \displaystyle \Theta (\omega(n)) & \text{\footnotesize se il lavoro nei nodi interni \(\omega(n)\) domina quello delle foglie \(\sigma(n)\)}
\end{cases}\]

\subsubsection*{Analisi del caso 1}
Si verifica quando il lavoro nei nodi interni \(\omega(n)\) è trascurabile rispetto a quello delle foglie \(\sigma(n)\):
\[\exists \varepsilon > 0 \text{ tale che } \omega(n) = O \! \left(n^{(\log_b a )- \varepsilon}\right) = \sigma(n) \cdot O \! \left(n^{-\varepsilon}\right)\]
Di seguito si osserva che calcolando il contributo complessivo dei nodi interni usando \(\omega(n) = n ^{(\log_b a) - \varepsilon}\),
il risultato è comparabile con quello delle foglie e non fa aumentare la complessità asintotica di \(\Theta(n^{\log_b a})\).
\begin{align*}
	T_{\text{int}}(n) &= \sum_{l=0}^{(\log_b n) - 1} a^l \cdot \omega \! \left( \frac{n}{b^l} \right) \; = \; \sum_{l=0}^{(\log_b n) - 1} a^l \cdot \left( \frac{n}{b^l} \right)^{(\log_b a) - \varepsilon} \; = \; \frac{n^{\log_b a}}{n^{\varepsilon}} \; \cdot \!\! \sum_{l=0}^{(\log_b n) - 1} \frac{a^l}{b^{l (\log_b a) - l \varepsilon}} \; = \\[3pt]
	&= \;\; \frac{n^{\log_b a}}{n^{\varepsilon}} \; \cdot \!\! \sum_{l=0}^{(\log_b n) - 1} \frac{a^l}{a^l \cdot b^{-l \varepsilon}} \; = \; \frac{n^{\log_b a}}{n^{\varepsilon}} \; \cdot \!\! \sum_{l=0}^{(\log_b n) - 1} {(b^{\varepsilon})}^l \; = \; \frac{n^{\log_b a}}{n^{\varepsilon}} \cdot \frac{(b^{\varepsilon})^{(\log_b n)} - 1}{b^{\varepsilon} - 1} \; = \\[3pt]
	&= \;\; \frac{n^{\log_b a}}{n^{\varepsilon}} \cdot \frac{n^{\varepsilon} - 1}{b^{\varepsilon} - 1} \; \approx \; \frac{n^{\log_b a}}{n^{\varepsilon}} \cdot \frac{n^{\varepsilon}}{b^{\varepsilon} - 1} \; = \; \frac{n^{\log_b a}}{b^{\varepsilon} - 1} \; = \; \frac{\sigma(n)}{b^{\varepsilon} - 1} \; = \; \Theta(\sigma(n))
\end{align*}
Si osserva che per \(\varepsilon \to 0^+\) il coefficiente moltiplicativo è molto elevato, ma comunque costante, per
cui asintoticamente trascurabile.

\subsubsection*{Analisi del caso 2}
Si verifica quando il lavoro nei nodi interni \(\omega(n)\) è dello stesso ordine di quello delle foglie \(\sigma(n)\):
\[\omega(n) = \Theta(n^{\log_b a}) = \Theta (\sigma(n))\]
Di seguito si dimostra che il contributo complessivo dei nodi interni ha una complessità asintotica maggiore di quella delle
foglie e domina la complessità totale dell'algoritmo, che risulta essere \(\Theta(n^{\log_b a} \, \log_b n)\).
\begin{align*}
	T_{\text{int}}(n) &= \sum_{l=0}^{(\log_b n) - 1} \left( a^l \cdot \omega \! \left( \frac{n}{b^l} \right) \right) \; = \; \sum_{l=0}^{(\log_b n) - 1} \left( a^l \cdot \left( \frac{n}{b^l} \right)^{\log_b a} \right) \; = \; n^{\log_b a} \; \cdot \!\! \sum_{l=0}^{(\log_b n) - 1} \left( \frac{a^l}{b^{l (\log_b a)}} \right) \\[3pt]
	&= n^{\log_b a} \; \cdot \!\! \sum_{l=0}^{(\log_b n) - 1} \frac{a^l}{a^l} \; = \; n^{\log_b a} \cdot (\log_b n) \; = \; \Theta(n^{\log_b a} \, \log_b n)
\end{align*}
Di seguito si osserva che il costo complessivo di ogni livello ha la stessa complessità asintotica di quello delle foglie,
per cui è possibile calcolare la complessità totale moltiplicando la complessità di ogni livello per il numero di livelli.
Gli algoritmi di questa tipologia con \(a = b\) (es. mergesort e quicksort) hanno la particolarità che il costo delle fasi di divide o
combine è proporzionale alla taglia dell'istanza.
\[T_{\text{\(l\)-esimo livello}}(n) = a^l \cdot \omega \! \left( \frac{n}{b^l} \right) = a^l \cdot \left( \frac{n}{b^l} \right)^{\log_b a} = a^l \cdot \frac{n^{\log_b a}}{b^{l (\log_b a)}} = a^l \cdot \frac{n^{\log_b a}}{a^l} = n^{\log_b a} = \sigma(n)\]

\newpage

\subsubsection*{Analisi del caso 3}
Si verifica quando il lavoro nei nodi interni \(\omega(n)\) domina quello delle foglie \(\sigma(n)\):
\[\exists \varepsilon > 0 \text{ tale che } \omega(n) = \Omega(n^{(\log_b a) + \varepsilon}) = \sigma(n) \cdot \Omega(n^{\varepsilon})\]
Dalla condizione precedente si può dedurre anche la seguente condizione di regolarità, ovvero che il lavoro interno complessivo
nei nodi figli (al livello inferiore) è minore del lavoro interno del nodo padre.
\[\exists c \in \mathbb{R} \text{ con } 0 < c < 1 \text{ tale che } \forall n \in \mathbb{N} \text{ si ha } a \cdot \omega \! \left( \frac{n}{b} \right) \leq c \cdot \omega(n)\]
Per la dimostrazione del terzo caso si utilizza la condizione di regolarità appena enunciata, in quanto rende lo svolgimento
più semplice e immediato. Di seguito sono riportati alcuni passaggi preliminari che verranno poi usati nella successiva
dimostrazione.
\[a\cdot \omega \! \left(\frac{n}{b}\right) \leq c\cdot \omega(n) \;\; \Rightarrow \;\; \omega \! \left(\frac{n}{b}\right) \leq \frac{c}{a} \cdot \omega(n) \;\; \Rightarrow \;\; \omega \! \left(\frac{n}{b^l}\right) \leq \left( \frac{c}{a} \right)^l \cdot \omega(n)\]
\[T_{\text{int}}(n) = \!\! \sum_{l=0}^{(\log_b n) - 1} a^l \cdot \omega \! \left( \frac{n}{b^l} \right) \; \leq \!\! \sum_{l=0}^{(\log_b n) - 1} a^l \cdot \frac{c^l}{a^l} \cdot \omega(n) \; = \; \omega(n) \cdot \!\!\! \sum_{l=0}^{(\log_b n) - 1} c^l \; < \; \omega(n) \cdot \frac{1}{1 - c} \; = \; \Theta(\omega(n))\]
Dall'analisi precedente si osserva che il costo complessivo dei nodi interni è proporzionale al costo di ogni nodo interno
e domina la complessità totale dell'algoritmo se il costo del singolo nodo interno è asintoticamente maggiore di quello delle
foglie, ovvero se \(\omega(n) = \Omega(n^{(\log_b a) + \varepsilon})\).

\subsubsection*{Considerazione sugli altri casi}
Esistono funzioni che non possono essere classificate in nessuno dei tre casi del master theorem, ad esempio quando
\(\omega(n) = n^{\log_b a} \log n\). In questi casi è necessario risolvere l'equazione alle ricorrenze attraverso
altri metodi risolutivi, ad esempio il formulone o l'analisi analitica.

\subsection{Applicazioni del divide and conquer}
Nei seguenti capitoli sono riportati alcuni esempi di applicazioni efficienti del paradigma divide and conquer:
\begin{itemize}
	\item esponenziazione con esponente intero
	\item moltiplicazione di polinomi tramite l'algoritmo di Karatsuba
	\item moltiplicazione di matrici tramite l'algoritmo di Strassen
	\item moltiplicazione di matrici supercircolanti tramite la trasformata di Hadamard
	\item moltiplicazione di polinomi tramite la trasformata di Fourier (FFT) e l'algoritmo di Cooley-Tukey
\end{itemize}

\newpage

\section{Esponenziazione con esponente intero}
\subsubsection*{Formalizzazione del problema}
Dato un elemento \(x\) proveniente da una struttura algebrica con operazione associativa (ad esempio numeri in \(\mathbb{N}\),
\(\mathbb{R}\), \(\mathbb{C}\), matrici, \dots) costante, si vuole calcolare \(x^n\) con \(n \in \mathbb{N}\). Siccome si
tratta \(x\) come costante, si esclude che possa essere un polinomio (in quanto aumenta di grado quando elevato a una potenza)
per semplificare la risoluzione.

\subsubsection*{Osservazioni sulla taglia del problema}
Essendo \(x\) costante, viene escluso dal calcolo della taglia che è determinata unicamente da \(n\). Si osserva che la
quantità di informazione necessaria a rappresentare \(n\) è \(\log_2 n\), per cui la taglia del problema in funzione di
cui va espressa la complessità è: \(\text{taglia}(n) = |n| = \log_2 n\), con \(n = 2 ^{\text{\, taglia}(n)}\).

\subsection{Algoritmo ovvio basato sulla definizione}
L'algoritmo ovvio per calcolare \(x^n\) si basa sulla definizione ricorsiva di potenza come composizione multipla
dell'operazione associativa della struttura algebrica, ovvero:
\[x^n := \begin{cases}
	\; x^{(0)} = 1 & \text{se } n = 0 \\[3pt]
	\; x^{(n-1)} \; \circ \; x & \text{se } n > 0 \quad - \quad \circ \;\; \text{è l'operazione associativa della struttura algebrica}
\end{cases}\]
La complessità è pari al costo di \(n - 1\) composizioni dell'operazione associativa, e risulta essere esponenziale in
funzione della taglia del problema, ovvero:
\[T(n) = n - 1 = 2^{\, \text{taglia}(n)} - 1 = \Theta \left( 2^{\, \text{taglia}(n)} \right)\]

\subsection{Algoritmo efficiente basato sul divide and conquer}
Si osserva che se l'operazione è associativa, come richiesto nelle ipotesi, è possibile riarrangiare le composizioni
di funzioni in modo da riutilizzare i risultati di composizioni già calcolate. Di seguito la definizione ricorsiva
dell'algoritmo efficiente basato sul divide and conquer:
\[x^n := \begin{cases}
	\; x^{(0)} = 1 & \text{se } n = 0 \\[3pt]
	\; x^{(1)} = x & \text{se } n = 1 \\[3pt]
	\; x^{(n/2)} \; \circ \; x^{(n/2)} & \text{se } n > 1 \text{ è pari} \\[3pt]
	\; x^{\lfloor n/2 \rfloor} \; \circ \; x^{\lfloor n/2 \rfloor} \; \circ \; x & \text{se } n > 1 \text{ è dispari}
\end{cases}\]
L'equazione alle ricorrenze per questo algoritmo al caso peggiore è la seguente, osservando che è sufficiente una chiamata
ricorsiva per calcolare \(x^{(n/2)}\) e, al caso peggiore, si aggiungono 2 operazioni di composizione.
\[T(n) = \begin{cases}
	\; 1 & \text{se } n = 0, 1 \\[3pt]
	\; T(n/2) + 2 & \text{se } n > 1
\end{cases}\]
Utilizzando il master theorem con \(a = 1\), \(b = 2\) e \(\omega(n) = O(1)\), si ottiene una complessità lineare
in funzione della taglia del problema, ovvero:
\[T(n) = \Theta \left( \log_2 n \right) = \Theta \left( \text{taglia}(n) \right)\]
